\section{Introduction}

Language-instructed agents increasingly translate task definitions, public
research mandates, and data into executable operations.  In quantitative
research, however, producing the next artifact is not the same as improving
the procedures that later research will use.  A stronger factor, a correct
pricing script, or a successful reproduction may be valuable while leaving
the agent no better able to formulate, implement, and assess the next study.
We investigate the second object: persistent improvement in how an autonomous
quantitative researcher works.

Research methodology is consequential even when the data and question are
held fixed.  In the \emph{Nonstandard Errors} experiment, 164 research teams
analyzed the same hypotheses and data but produced substantial dispersion
through different analysis paths \citep{menkveld2024nonstandard}.  This does
not make every disagreement an error.  It shows that measure construction,
sampling, observation treatment, and statistical specification are part of
the evidence-generating process.  Investment-research protocols likewise
emphasize choices about economic motivation, data, statistical analysis, and
evaluation \citep{arnott2019backtesting}.  Recent work finds analogous
methodological variation among autonomous coding agents performing financial
analysis \citep{gao2026nonstandard}.

We ask: \emph{can a quantitative research agent improve its reusable research
capabilities by experimenting on the way it conducts research?}  We model the
researcher as a system with two roles.  A \Worker{} carries out a quantitative
assignment.  An \Evolver{} examines selected experience and develops changes
to how later \Worker{} instances operate.  The learned object is the
\fullharness{}: prompts, tools, bindings, executable and textual skills,
memory/context operations, middleware, validators, routing, and control flow.
The foundation-model configurations and revision policy are fixed within a
declared experiment or final campaign. Thus
``self-improvement'' refers to changing system capabilities, not to recursive
optimization of the \Evolver{} algorithm or model weights.

We study \method{} within full-harness adaptation: whether an agent can learn
reusable research operations whose use depends on the conditions of the
quantitative question. An operation connects available inputs to a computation
or decision and a required output; its applicability concerns whether the target
quantity, sample, information convention, and intended use support that operation.
The proposed \Evolver{} uses public evidence and its own observed work to decide
whether to reuse, adapt, replace, or decline an operation, and tests a persistent
revision through a fresh \Worker{}. The empirical question is whether these
decisions improve later work, not whether the system records applicability
fields or calls a checker. Ordinary interface and delivery repairs remain
legitimate, and no experiment, code edit, or memory update is mandatory in every
round. Investigator-selected tasks and evidence are attributed separately from
the \Evolver{}'s subsequent diagnosis and revision. Feedback retains only the
episode layers actually produced; missing revisions, Workers, and outcomes
remain explicit.

Our intended contributions are:

\begin{enumerate}
  \item \textbf{A research-capability learning problem.}  The evolving object
  is a complete task-conditioned researcher, not a collection of best answers.
  \item \textbf{A testable quantitative revision hypothesis.} We investigate
  whether the target quantity, sample, and admissible information change an
  executable comparison or justified reuse decision, with consequences for
  later \Worker{} behavior.
  \item \textbf{An end-to-end evidence design.}  We jointly require matched
  performance, real persistent changes with fresh \Worker{} use, and actual
  cumulative multi-task operation, while retaining unsuccessful episodes and
  cost.
  \item \textbf{A demonstration on backtesting tasks.}  The evolved harness
  scores 16/20 against the initial harness's 14/20 on a controlled
  BacktestBench comparison, with two gains and zero regressions.
\end{enumerate}

The central empirical result is a complete controlled BacktestBench
comparison: the evolved harness scores 16/20 against the initial harness's
14/20, with two task gains and zero regressions across 20 new instances
(Section~\ref{sec:current-results}).  The \Evolver{} diagnosed that the
initial \Worker{} assessed round-trip performance using full-window daily
mark-to-market P\&L, which conflates unrealized and realized gains across
multi-day holds; the accepted revision scopes evaluation to window-only
indicators and realized round-trip PNL.  Fresh-\Worker{} consumer checks
confirm consequential policy uptake in four cases.  This demonstrates that
research-method experimentation can identify a measurement-methodology
limitation and produce a matched initial-to-frozen improvement on backtesting
tasks.  QuantCodeEval results across the same evolved harness are mixed, which
is consistent with the conditional nature of research-method improvement: a
backtesting measurement policy does not transfer to tasks requiring different
research operations.  Executable harness components beyond prompt-level
changes and cumulative multi-task benefit remain open empirical targets.
We organize findings around performance, learned operations, and cumulative
research, retaining all negative results and the full development chronology
in the appendix.
